\section{Problem setup and the shared pipeline}
\label{sec:setup}

\paragraph{Grid and bricks.} A brick model lives on an integer grid. The horizontal unit is one \emph{stud}
(8\,mm pitch); the vertical unit is one \emph{plate} (3.2\,mm), and a standard brick is three plates tall. We use a
16\,$\times$\,16 stud footprint and 8 brick layers, which is 2048 cells of one stud by one stud by one brick. The
catalogue has six upright bricks: 1\,$\times$\,1, 1\,$\times$\,2, 1\,$\times$\,3, 1\,$\times$\,4, 2\,$\times$\,2 and
2\,$\times$\,4, placed at 0\textdegree{} or 90\textdegree. A \emph{placement} is a brick type, rotation and
position; a \emph{build} is a list of placements plus an ordered list of assembly steps.

\paragraph{Target.} The input is a watertight mesh. It is rotated so its chosen axis points up, scaled uniformly to fit
the physical box the grid represents, floored, and voxelized by testing whether each cell centre lies inside the
mesh. The result is a boolean \emph{target} grid $T$. (Procedural targets are generated directly on the grid.)

\paragraph{Legality.} A placement is legal when it lies inside the grid, overlaps no placed brick, is \emph{supported}
(it sits on the ground or at least one cell directly beneath it is already filled), covers at least one target cell,
and is at most half \emph{filler}, cells outside the target. These rules are enforced by a candidate enumerator that
every generator in this paper shares, so no generator can place an illegal brick.

\paragraph{The finalize step.} Every generator returns a raw list of placements, and every list goes through the
same four stages (Figure~\ref{fig:pipeline}):
\begin{enumerate}
  \item \textbf{Trim.} Bricks connect only through studs: an upper brick is attached to a lower one when it rests
  directly on it with overlapping footprints. Side contact does not count. The build is reduced to its largest
  component connected to the ground; everything else is discarded. This is what makes the output one object.
  \item \textbf{Plan.} A planner orders the remaining bricks one per step, each inserted straight down, attached to
  what is already there, with a clear insertion path and never sealing an unplaced brick beneath it. For this
  catalogue of same-height upright bricks that greedy planner is complete: every valid build gets an order.
  \item \textbf{Validate.} An independent validator replays the steps and checks bounds, collisions, support, the
  single-component property and the piece budget. A build that fails here is reported as invalid.
  \item \textbf{Stability check.} A linear simplification of the StableLego force model~\citep{liu2024stablelego}: bricks
  are rigid, loaded by their own weight; each shared stud cell carries vertical force at four points; compression is
  unbounded and tension is capped by a clutch capacity of 0.98\,N per stud cell, divided over its four points. The
  constant is StableLego's measured value, which that work applies per contact point rather than per stud, so our
  limit is roughly four times more conservative than theirs. A linear program
  finds the smallest tension utilisation that balances every brick. Utilisation at most one is \emph{stable under the
  model}. The model ignores lateral loads, friction and deformation, so throughout this paper ``stable'' means
  \stablenote.
\end{enumerate}

\begin{figure}[t]
  \centering
  \resizebox{\linewidth}{!}{%
  \begin{tikzpicture}[
    box/.style={draw=muted, fill=boxc, rounded corners=2pt, align=center, font=\small, minimum height=9mm, inner sep=4pt},
    gen/.style={box, draw=rlc, line width=0.9pt},
    arr/.style={-{Stealth[length=2mm]}, muted, line width=0.7pt},
    lab/.style={font=\scriptsize, text=muted, align=center}]
    \node[box] (mesh) at (0,0) {mesh or\\procedural shape};
    \node[box] (vox) at (3.1,0) {voxelize\\$16\times16\times8$};
    \node[gen] (gen) at (6.6,0) {\textbf{generator}\\greedy search / SFT policy /\\RL policy};
    \node[box] (trim) at (10.4,0) {trim to largest\\ground-connected\\component};
    \node[box] (plan) at (13.4,0) {plan\\assembly order};
    \node[box] (val) at (13.4,-2.6) {validate: bounds, collisions,\\support, one component, budget};
    \node[box] (stab) at (8.6,-2.6) {stability LP\\(simplified static force model)};
    \node[box, draw=greedyc] (out) at (3.4,-2.6) {build + parts list\\+ checked steps + metrics};
    \draw[arr] (mesh) -- (vox);
    \draw[arr] (vox) -- (gen);
    \draw[arr] (gen) -- (trim);
    \draw[arr] (trim) -- (plan);
    \draw[arr] (plan) -- (val);
    \draw[arr] (val) -- (stab);
    \draw[arr] (stab) -- (out);
    \node[lab, above=1mm of gen] {the only part that differs between arms};
    \node[lab, above=1mm of trim, xshift=1.5cm] {shared by every arm};
    \node[lab, below=1mm of gen, text=rlc] {chooses only among legal candidates: inside the grid,\\no overlap, supported, $\geq 1$ target cell, $\leq$ half filler};
  \end{tikzpicture}}
  \caption{The shared pipeline. Generators differ only in how they choose the next brick; the trim, planner,
  validator and stability check are identical, so fidelity differences cannot come from different rules.}
  \label{fig:pipeline}
\end{figure}

\paragraph{Metrics.} \emph{Fidelity} is the intersection-over-union between the occupancy of the trimmed build $B$ and
the target, $\iou = |B \cap T| / |B \cup T|$, so both missing and extraneous cells count. \emph{Validity} is the
validator's verdict, \emph{stability} the force model's, \emph{pieces} the brick count after the trim, and
\emph{decode time} the wall-clock seconds the generator took. The \emph{piece budget} is a hard cap on placements;
a generator may stop earlier.

\paragraph{Fidelity convention.} Throughout, a reported mean IoU is the mean over the \nAttempted{} benchmark cases
that are attemptable (one fixture is a deliberately non-watertight mesh that every generator must reject), with a
build that is missing or fails validation scored 0. A generator therefore cannot raise its mean by failing on hard
cases. The benchmark's own summariser averages over valid builds only, which reads higher for generators with
failures; we never use that number for comparisons.
